\documentclass[12pt,a4paper]{article}
\usepackage[dvipsnames]{xcolor}
\usepackage[utf8]{inputenc}
\usepackage{amsmath}
\usepackage{mathtools}
\usepackage{marvosym}
\usepackage{wrapfig}
\usepackage{hyperref}
\usepackage{float}
\usepackage{multicol}
\hypersetup{colorlinks,citecolor=black,filecolor=black,linkcolor=black,urlcolor=black}
\usepackage{pdfpages}
\usepackage{amsfonts}
\usepackage{amssymb}
\usepackage{fancyhdr}
\usepackage{chemfig}
\usepackage{graphicx}
\usepackage[T1]{fontenc}
\usepackage[magyar]{babel}
\usepackage{bm}
\usepackage{tikz, tcolorbox}
\usetikzlibrary{positioning}
\usepackage{verbatim}
\usepackage{amsmath,esint}
\usepackage{setspace}
\usepackage{qtree}
\usepackage{multido}
\usepackage{mathptmx} % Times New Roman betűtípus
\usepackage{listings}
\usepackage{caption}
\usepackage{xparse}
\captionsetup[lstlisting]{labelformat=empty}

\lstset{
    language=Matlab,
    inputencoding=utf8,
    extendedchars=true,
    basicstyle=\ttfamily\small,
    frame=single,
    breaklines=true,
    numbers=left,
    numberstyle=\tiny,
    literate=
        {á}{{\'a}}1
        {é}{{\'e}}1
        {í}{{\'i}}1
        {ó}{{\'o}}1
        {ö}{{\"o}}1
        {ő}{{\H{o}}}1
        {ú}{{\'u}}1
        {ü}{{\"u}}1
        {ű}{{\H{u}}}1
        {Á}{{\'A}}1
        {É}{{\'E}}1
        {Í}{{\'I}}1
        {Ó}{{\'O}}1
        {Ö}{{\"O}}1
        {Ő}{{\H{O}}}1
        {Ú}{{\'U}}1
        {Ü}{{\"U}}1
        {Ű}{{\H{U}}}1
}
\renewcommand{\figureautorefname}{ábra}
\newcommand{\abrara}[1]{\hyperref[#1]{\ref*{#1}.~ábra}}
\NewDocumentCommand{\kep}{O{300px} O{0} m m m}{%
\begin{figure}[H]
    \centering
    \includegraphics[width=#1, angle=#2]{#3}
    \caption{#4}
    \label{fig:#5}
\end{figure}%
}
\newcommand{\Pointilles}[1]{%
\par\nobreak
\noindent\rule{0pt}{1.5\baselineskip}% Provides a larger gap between the
\multido{}{#1}{\noindent\makebox[\linewidth]{\dotfill}\endgraf}
\bigskip% Gap between dots and next paragraph
}
\newcommand{\feladat}[1]{%
\phantomsection
\subsection*{#1. feladat}%
\addcontentsline{toc}{subsection}{#1. feladat}%
}
% Kitöltendő részek jelölése (leadás előtt ki kell venni!)
\newcommand{\todo}[1]{\textcolor{red}{\textbf{TODO:} #1}}
\usepackage{pgfplots}
\pgfplotsset{height = 10cm, width=15cm,compat=1.9}
\usepackage[left=2cm,right=2cm,top=2cm,bottom=2cm]{geometry}
\setlength{\parindent}{0pt}
\setlength{\parskip}{0em}
\pagestyle{fancy}
\fancyhf{}
\cfoot{\thepage. oldal}

% Másfeles sorköz beállítása
\onehalfspacing

\begin{document}

\begin{titlepage}
\centering

\includegraphics[width=.8\textwidth]{../../Kepek/bme_logo_nagy.jpg}

\vspace{1em}

{
\large
BUDAPESTI MŰSZAKI ÉS GAZDASÁGTUDOMÁNYI EGYETEM\\ VILLAMOSMÉRNÖKI ÉS INFORMATIKAI KAR\\ MESTERSÉGES INTELLIGENCIA ÉS RENDSZERTERVEZÉS TANSZÉK
}

\vspace{6em}

{
\Huge
\textbf{Adaptív szűrők vizsgálata}
}

\vspace{3em}

{
\huge
Intelligens beágyazott rendszerek laboratórium\\

\vspace{.2em}

\large
(BMEVIMIMA21)
}

\vspace{9em}

{
\Large
\textit{Készítette:}\\
\vspace{.5em}
\textbf{Pál Patrik, Gavallér Dávid}\\
\vspace{.5em}
\textbf{AIQTQW, ES02RO}
}

\vspace{12em}

{
\large
BUDAPEST, \today
}
\end{titlepage}

\tableofcontents

\newpage

\section{Bevezetés}

\subsection*{A mérés célja}

A mérés célja a digitális jelfeldolgozás témakörébe tartozó adaptív szűrőkkel kapcsolatos ismeretek elmélyítése és azok gyakorlati vonatkozásainak megismerése. Az adaptív szűrők segítségével egy ismeretlen, akár időben változó átvitelű rendszer modellje állítható elő úgy, hogy a modell paramétereit a valós rendszer és a modell kimenetének eltérése alapján folyamatosan hangoljuk.\par\medskip

Az első feladatrészben MATLAB szimulációk segítségével megvalósítjuk és megvizsgáljuk az LMS, az NLMS, valamint az FxLMS algoritmust. Az algoritmusokkal egy előre megtervezett IIR szűrőt identifikálunk, és megvizsgáljuk, hogyan befolyásolja a bátorsági tényező, a szűrőhossz és a gerjesztőjel teljesítménye a konvergenciát és az approximáció jóságát.\par\medskip

A második feladatrészben egy ADSP-21364 jelfeldolgozó processzoron futó, előre elkészített programok segítségével, egy elektronikus hibrid áramkörön adaptív vonaljavítást, adaptív visszhangcsökkentést, valamint aktív zajcsökkentést valósítunk meg.

\subsection*{A méréshez használt eszközök}
\begin{itemize}
	\item Számítógép
	\item MATLAB
	\item CrossCore Embedded Studio
	\item ADSP-21364 EZ-KIT Lite (SHARC) fejlesztőkártya
	\item Elektronikus hibrid áramkör
	\item Függvény- és zajgenerátor
	\item Audio jelforrás
	\item Oszcilloszkóp
	\item Fejhallgató
\end{itemize}

\todo{Az eszközlistát a mérés után pontosítani kell.}

\newpage

\section{Mérési feladatok}

%==========================================================================
\feladat{1}

Az első feladat célja az LMS algoritmus megvalósítása MATLAB-ban és annak vizsgálata egy rendszeridentifikációs feladaton keresztül. Az LMS algoritmus a Wiener-szűrő optimális együtthatóit a legmeredekebb lejtő módszerével, a statisztikai jellemzők (autokorrelációs mátrix, keresztkorrelációs vektor) ismerete nélkül, a pillanatnyi hiba alapján közelíti.

\feladat{1.1}

\subsubsection*{Feladat leírása}

Az elméleti áttekintésben szereplő (13) és (14) egyenleteknek megfelelően írjon MATLAB függvényt az LMS algoritmus rendszeridentifikációs feladatot ellátó megvalósítására! Az algoritmust megvalósító függvény váza:\\
\texttt{function [e,w]=lms(mu,M,x,y)}

\subsubsection*{Feladat megoldása}

Az \texttt{lms.m} függvényben az adaptív FIR szűrő $\mathbf{x}(n) = [x(n),\, x(n-1),\, \ldots,\, x(n-M+1)]^T$ regressziós vektorát minden lépésben eltoljuk, és a legfrissebb bemeneti mintát a vektor elejére helyezzük. Ezt követően kiszámítjuk a becslési hibát, majd frissítjük a szűrőegyütthatókat az alábbi összefüggések szerint:

\begin{equation}
    e(n) = y(n) - \hat{\mathbf{w}}^T(n)\,\mathbf{x}(n),
\end{equation}

\begin{equation}
    \hat{\mathbf{w}}(n+1) = \hat{\mathbf{w}}(n) + 2\mu\, e(n)\,\mathbf{x}(n).
\end{equation}

A szűrőegyütthatókat nulláról indítjuk. A függvény a teljes hibajelet és a szűrőegyütthatók végső értékét adja vissza. A függvény kódja a Függelékben található.

\feladat{1.2}

\subsubsection*{Feladat leírása}

Tervezzen IIR szűrőt a mérésvezető utasításai szerint!

\subsubsection*{Feladat megoldása}

\todo{A mérésvezető által megadott paramétereket beírni, és a \texttt{tervezIIR.m} függvényben átállítani!}\par\medskip

Az IIR szűrőt a \texttt{tervezIIR.m} függvényben terveztük meg, így a későbbi feladatokban is mindig ugyanazt a szűrőt használtuk. A szűrő ideiglenesen egy 48~kHz-es mintavételi frekvenciához tervezett, negyedrendű, 4~kHz határfrekvenciájú Butterworth aluláteresztő szűrő, amit a MATLAB \texttt{butter} függvényével állítottunk elő. A szűrő amplitúdómenete a(z) \abrara{fig:iir-atvitel}-n látható.

\kep[0.8\textwidth]{../../Kepek/IIR_atvitel.jpg}{A tervezett IIR szűrő amplitúdómenete.}{iir-atvitel}

\feladat{1.3}

\subsubsection*{Feladat leírása}

Az IIR szűrőt az LMS algoritmus segítségével identifikálja! Gerjesztésként használjon Gauss-\hspace{0pt}eloszlású fehér zajt! Vesse össze az IIR szűrő és az identifikáció eredményeként kiadódó FIR szűrő átvitelét!

\subsubsection*{Feladat megoldása}

Gerjesztésként a \texttt{randn} függvénnyel egységnyi szórású, $N = 20000$ mintából álló Gauss-eloszlású fehér zajt állítottunk elő, amit a \texttt{filter} függvénnyel átvezettünk a tervezett IIR szűrőn. A gerjesztést és az IIR szűrő kimenetét az \texttt{lms} függvény bemenetére adtuk $M = 32$ együtthatóval és $\mu = 0{,}005$ bátorsági tényezővel.\par\medskip

Az algoritmus konvergenciáját a négyzetes hiba mozgóátlagával (tanulási görbe) jellemeztük, ami a(z) \abrara{fig:lms-tanulasi-gorbe}-n látható.

\kep[0.8\textwidth]{../../Kepek/LMS_tanulasi_gorbe.jpg}{Az LMS algoritmus tanulási görbéje.}{lms-tanulasi-gorbe}

Az IIR szűrő és az identifikált FIR szűrő átvitelét a(z) \abrara{fig:lms-atvitel}-n, az IIR szűrő impulzusválaszát és az identifikált FIR szűrő együtthatóit pedig a(z) \abrara{fig:lms-impulzus}-n hasonlítottuk össze.

\kep[0.8\textwidth]{../../Kepek/LMS_atvitel_osszehasonlitas.jpg}{Az IIR szűrő és az identifikált FIR szűrő átvitele.}{lms-atvitel}

\kep[0.8\textwidth]{../../Kepek/LMS_impulzusvalasz.jpg}{Az IIR szűrő impulzusválasza és az identifikált FIR szűrő együtthatói.}{lms-impulzus}

\todo{Az eredmények kiértékelése (állandósult hiba, átviteli hiba, a záró sávban látható eltérés oka).}

\feladat{1.4}

\subsubsection*{Feladat leírása}

Próbáljon ki különféle $\mu$ és $M$ értékeket, és vizsgálja meg, hogyan befolyásolják az adaptív szűrő konvergenciáját, illetve az IIR szűrő approximációjának jóságát!

\subsubsection*{Feladat megoldása}

Először rögzített $M = 32$ szűrőhossz mellett különböző bátorsági tényezőkkel futtattuk le az identifikációt. A tanulási görbéket a(z) \abrara{fig:lms-mu}-n ábrázoltuk.

\kep[0.8\textwidth]{../../Kepek/LMS_mu_hatasa.jpg}{A bátorsági tényező hatása az LMS algoritmus konvergenciájára.}{lms-mu}

\todo{Kiértékelés: a nagyobb $\mu$ gyorsabb beállást, de nagyobb állandósult hibát (misadjustment) okoz, túl nagy $\mu$ esetén az algoritmus instabil.}\par\medskip

Ezt követően rögzített $\mu = 0{,}002$ bátorsági tényező mellett a szűrőhosszat változtattuk. A tanulási görbék a(z) \abrara{fig:lms-m}-n, az identifikált FIR szűrők átvitele pedig a(z) \abrara{fig:lms-m-atvitel}-n látható.

\kep[0.8\textwidth]{../../Kepek/LMS_M_hatasa.jpg}{A szűrőhossz hatása az LMS algoritmus konvergenciájára.}{lms-m}

\kep[0.8\textwidth]{../../Kepek/LMS_M_atvitel.jpg}{Különböző hosszúságú identifikált FIR szűrők átvitele.}{lms-m-atvitel}

\todo{Kiértékelés: rövid szűrő nem képes az IIR szűrő impulzusválaszát leírni (nagy modellhiba), hosszú szűrőnél kisebb a modellhiba, de lassabb a beállás és nagyobb a gradienszaj.}

\newpage

%==========================================================================
\feladat{2}

A második feladatban az NLMS algoritmust valósítottuk meg és hasonlítottuk össze az LMS algoritmussal. Az NLMS a bátorsági tényezőt a bemeneti jel energiájával normálja, így a konvergencia sebessége független a gerjesztés amplitúdójától.

\feladat{2.1}

\subsubsection*{Feladat leírása}

Az elméleti áttekintésben szereplő (15) és (16) egyenleteknek megfelelően írjon MATLAB függvényt az NLMS algoritmus rendszeridentifikációs feladatot ellátó megvalósítására! Az algoritmust megvalósító függvény váza:\\
\texttt{function [e,w]=nlms(mu,a,M,x,y)}

\subsubsection*{Feladat megoldása}

Az \texttt{nlms.m} függvény felépítése megegyezik az \texttt{lms.m} függvényével, csak az együtthatók frissítése tér el:

\begin{equation}
    e(n) = y(n) - \hat{\mathbf{w}}^T(n)\,\mathbf{x}(n),
\end{equation}

\begin{equation}
    \hat{\mathbf{w}}(n+1) = \hat{\mathbf{w}}(n) + \frac{\tilde{\mu}}{a + \mathbf{x}^T(n)\,\mathbf{x}(n)}\, e(n)\,\mathbf{x}(n).
\end{equation}

Az $a$ regularizációs konstans megakadályozza, hogy kis bemeneti jel esetén a nevező nullához közeli értékű, így a lépésköz túl nagy legyen. A MATLAB függvényben a $\tilde{\mu}$ lépésközt az előírt függvényváznak megfelelően a \texttt{mu} paraméter jelöli.

\feladat{2.2}

\subsubsection*{Feladat leírása}

Próbáljon ki különféle gerjesztőjel-teljesítményt, illetve $\mu$ értékeket, és vizsgálja meg, hogyan befolyásolják az adaptív szűrő konvergenciáját! Az előző feladatban tervezett IIR szűrőt identifikálja, és az eredményeket vesse össze az ott kapottakkal!

\subsubsection*{Feladat megoldása}

Az 1.3 feladatban használt fehér zajt különböző szórásokkal skáláztuk, majd mind az LMS, mind az NLMS algoritmussal elvégeztük az identifikációt. Az összehasonlíthatóság érdekében a négyzetes hibát az IIR szűrő kimenetének teljesítményére normáltuk. Az eredmény a(z) \abrara{fig:nlms-teljesitmeny}-n látható.

\kep[0.9\textwidth]{../../Kepek/NLMS_teljesitmeny.jpg}{A gerjesztőjel teljesítményének hatása az LMS és az NLMS algoritmus konvergenciájára.}{nlms-teljesitmeny}

\todo{Kiértékelés: az LMS beállási sebessége a gerjesztés teljesítményétől függ, az NLMS-é nem.}\par\medskip

Az NLMS algoritmus esetén is megvizsgáltuk a bátorsági tényező hatását, ami a(z) \abrara{fig:nlms-mu}-n látható.

\kep[0.8\textwidth]{../../Kepek/NLMS_mu_hatasa.jpg}{A bátorsági tényező hatása az NLMS algoritmus konvergenciájára.}{nlms-mu}

\todo{Kiértékelés: stabilitási tartomány $0 < \tilde{\mu} < 2$, a beállási sebesség és az állandósult hiba közötti kompromisszum.}\par\medskip

Végül az LMS és az NLMS algoritmussal identifikált FIR szűrők átvitelét vetettük össze az IIR szűrő átvitelével, ami a(z) \abrara{fig:nlms-lms}-n látható.

\kep[0.8\textwidth]{../../Kepek/NLMS_LMS_osszehasonlitas.jpg}{Az LMS és az NLMS algoritmussal identifikált FIR szűrők átvitele.}{nlms-lms}

\todo{Kiértékelés, összevetés az 1.3 feladat eredményeivel.}

\newpage

%==========================================================================
\feladat{3}

\textit{(kiegészítő mérési feladat)} A harmadik feladatban az FxLMS algoritmust valósítottuk meg és vizsgáltuk. Az FxLMS algoritmust olyan esetekben kell alkalmazni, amikor az adaptív szűrő kimenete és a hibajel képzése közé egy nem egységnyi átvitelű, ún. másodlagos út ($S(z)$) iktatódik be, ami például aktív zajcsökkentés esetén a beavatkozó hangszóró és a hibamikrofon közötti átvitel.

\feladat{3.1}

\subsubsection*{Feladat leírása}

Az elméleti áttekintésben szereplő (17), (18), (19) és (20) egyenleteknek megfelelően írjon MATLAB függvényt az FxLMS algoritmus megvalósítására! Az $S(z)$ átvitel egy erre a feladatra tervezett egyszerű FIR szűrő legyen! Az algoritmust megvalósító függvény váza:\\
\texttt{function [e,w]=fxlms(mu,M,s,c,x,y)}

\subsubsection*{Feladat megoldása}

Mivel az $S(z)$ másodlagos út is része az adaptációs huroknak, az \texttt{fxlms.m} függvényben a számítást mintánként végeztük. Minden lépésben kiszámítottuk az adaptív szűrő $\hat{y}(n)$ kimenetét, amit átvezettünk az $S(z)$ másodlagos úton, így kaptuk a $\hat{y}_s(n)$ jelet. A hibajel:

\begin{equation}
    e(n) = y(n) - \hat{y}_s(n).
\end{equation}

A referenciajelet a másodlagos út $C(z)$ becslőjével szűrtük, és a szűrt referenciajel utolsó $M$ mintájából képeztük az $\mathbf{x}_c(n)$ vektort:

\begin{equation}
    \mathbf{x}_c(n) = \left[\sum_{i=0}^{I-1} c_i x(n-i),\; \sum_{i=0}^{I-1} c_i x(n-i-1),\; \ldots,\; \sum_{i=0}^{I-1} c_i x(n-i-(M-1))\right]^T.
\end{equation}

Az együtthatók frissítése:

\begin{equation}
    \hat{\mathbf{w}}(n+1) = \hat{\mathbf{w}}(n) + \mu\, e(n)\,\mathbf{x}_c(n).
\end{equation}

\feladat{3.2}

\subsubsection*{Feladat leírása}

Identifikálja az $S(z)$ átvitelt az előzőleg megírt LMS algoritmussal, az identifikált együtthatókat használja fel a $C(z)$ szűrőben!

\subsubsection*{Feladat megoldása}

Az $S(z)$ másodlagos utat egy 5 mintás késleltetésből és egy tizedrendű, \texttt{fir1} függvénnyel tervezett aluláteresztő FIR szűrőből állítottuk össze, ezzel modellezve az AD- és DA-átalakítók késleltetését és sávkorlátozását. Az $S(z)$ átvitelt fehér zaj gerjesztéssel, az \texttt{lms} függvénnyel identifikáltuk $M = 24$ együtthatóval, és az eredményt a $C(z)$ szűrő együtthatóiként használtuk fel. Az identifikáció eredménye a(z) \abrara{fig:fxlms-masodlagos}-n látható.

\kep[0.8\textwidth]{../../Kepek/FxLMS_masodlagos_ut.jpg}{A másodlagos út és az identifikált modelljének együtthatói és átvitele.}{fxlms-masodlagos}

\feladat{3.3}

\subsubsection*{Feladat leírása}

Tesztelje az FxLMS algoritmust többféle referenciajellel (szinuszjel, háromszögjel stb.)! A $P(z)$ átvitel legyen az előzőekben használt IIR szűrő!

\subsubsection*{Feladat megoldása}

A $P(z)$ elsődleges út az 1.2 feladatban tervezett IIR szűrő volt. Referenciajelként 500~Hz-es szinusz-, háromszög- és négyszögjelet használtunk, az elnyomandó jelet a referenciajel $P(z)$-n való átvezetésével állítottuk elő. Az algoritmust $M = 64$ együtthatóval és $\mu = 0{,}0005$ bátorsági tényezővel futtattuk. Az elnyomandó jel és a hibajel a(z) \abrara{fig:fxlms-referencia}-n látható.

\kep[0.9\textwidth]{../../Kepek/FxLMS_referenciajelek.jpg}{Az FxLMS algoritmus működése különböző referenciajelekkel.}{fxlms-referencia}

\todo{Kiértékelés, az elnyomás mértéke az egyes referenciajelekre.}

\feladat{3.4}

\subsubsection*{Feladat leírása}

Készítsen szimulációt, amely demonstrálja, hogy $S(z)$ elégtelen identifikációja esetén az algoritmus nem konvergens, még $\mu$ minden határon túl történő csökkentése esetén sem!

\subsubsection*{Feladat megoldása}

Szinuszos referenciajel esetén az FxLMS algoritmus akkor konvergens, ha a referenciajel frekvenciáján az $S(z)$ másodlagos út és annak $C(z)$ modellje közötti fáziseltérés kisebb $90^\circ$-nál. Az elégtelen identifikációt úgy modelleztük, hogy a 3.2 feladatban kapott $C(z)$ modellt 20 mintával tovább késleltettük, így az 1~kHz-es referenciajel frekvenciáján a fáziseltérés meghaladja a $90^\circ$-ot, ahogy azt a(z) \abrara{fig:fxlms-fazis} mutatja.

\kep[0.8\textwidth]{../../Kepek/FxLMS_faziselteres.jpg}{Az $S(z)$ másodlagos út és a hibás $C(z)$ modell fáziseltérése.}{fxlms-fazis}

Az algoritmust a hibás modellel több, egyre kisebb bátorsági tényezővel is lefuttattuk, és összevetettük a helyes modellel kapott eredménnyel. Az eredmény a(z) \abrara{fig:fxlms-divergencia}-n látható.

\kep[0.8\textwidth]{../../Kepek/FxLMS_divergencia.jpg}{Az FxLMS algoritmus divergenciája hibás másodlagosút-modell esetén.}{fxlms-divergencia}

\todo{Kiértékelés: $\mu$ csökkentése csak a divergencia sebességét csökkenti, a hiba minden esetben korlátlanul nő.}

\newpage

%==========================================================================
\feladat{4}

A negyedik feladattól kezdve az ADSP-21364 EZ-KIT Lite fejlesztőkártyával és az elektronikus hibrid áramkörrel végeztünk kísérleteket. A kártyára az előre elkészített projekteket a CrossCore Embedded Studio segítségével töltöttük le, a paramétereket pedig a memória tartalmának módosításával állítottuk.\par\medskip

A negyedik feladatban az adaptív vonaljavító (ALE) struktúrát vizsgáltuk. A struktúra a bemeneti jel $\Delta$ mintával késleltetett változatából egy lineáris prediktor segítségével becsüli a bemeneti jelet, amivel a hasznos jel és a zaj szétválasztható.

\feladat{4.1}

\subsubsection*{Feladat leírása}

A hibrid segítségével összegezze a hasznos jelet és a zajt! Ügyeljen arra, hogy se az áramkört, se a forrást ne terhelje jelentősen!

\subsubsection*{Feladat megoldása}

A DSP kártyára a \texttt{SHARC\_Lab\_ALE} projektet töltöttük le. A projektben a kodek 48~kHz-es mintavételezése után 6-szoros decimálás történik (\texttt{INTERP\_RAT = 6} a \texttt{tt.h} állományban), így a feldolgozás tényleges mintavételi frekvenciája $f_s = 8~\mathrm{kHz}$. Az ALE struktúra $y(n)$ bemeneti jelét a kártya bemenetének bal (piros) csatornájára kötöttük. A kártya kimenetein a következő jelek érhetők el (a 4. kimenetre a \texttt{FLG1} ($x$), \texttt{FLG2} ($e$), \texttt{DAI\_P19} ($\hat{y}$) és \texttt{DAI\_P20} ($y$) gombokkal is átkapcsolható a kívánt jel):

\begin{itemize}
    \item $y(n)$: 1. kimenet bal és jobb csatornája,
    \item $\hat{y}(n)$: 2. kimenet bal és jobb csatornája,
    \item $e(n)$: 3. kimenet bal és jobb csatornája,
    \item $x(n)$: 4. kimenet bal és jobb csatornája.
\end{itemize}

\todo{A hasznos jel és a zaj összegzésének módja a hibriden, a mérési elrendezés leírása.}

% \kep{../../Kepek/XXXXXXXX_XXXXXX.jpg}{A mérési elrendezés.}{meresi-elrendezes-ale}

\feladat{4.2}

\subsubsection*{Feladat leírása}

Periodikus zajt alkalmazva működtesse a struktúrát! Vizsgálja meg, milyen $\Delta$ késleltetésre van szükség!

\subsubsection*{Feladat megoldása}

\todo{Mérési eredmények. A projektben a késleltetés a \texttt{D} (alapértelmezetten 1), a szűrőhossz az \texttt{N} (alapértelmezetten 100) konstanssal állítható a \texttt{processALE.c} fájlban. Keskenysávú zaj esetén a kimenet az $e(n)$ jel.}

\feladat{4.3}

\subsubsection*{Feladat leírása}

Véletlenszerű zajt alkalmazva működtesse a struktúrát! Vizsgálja meg, milyen $\Delta$ késleltetésre van szükség!

\subsubsection*{Feladat megoldása}

\todo{Mérési eredmények. Szélessávú zaj esetén a kimenet a $\hat{y}(n)$ jel.}

\feladat{4.4}

\subsubsection*{Feladat leírása}

Próbáljon ki különféle $\mu$ értékeket, és vizsgálja a rendszer beállási és stabilitási tulajdonságait! Milyen hatással van $\mu$ változtatása a hasznos jelre?

\subsubsection*{Feladat megoldása}

\todo{Mérési eredmények. A bátorsági tényező a \texttt{mu} változóval állítható (alapértelmezetten 0,01).}

\newpage

%==========================================================================
\feladat{5}

Az ötödik feladatban az adaptív visszhangcsökkentést vizsgáltuk. A visszhangcsökkentés a hallgató oldalán elhelyezkedő hibrid átvitelének rendszeridentifikációs feladata, amit a beszélő oldalán elérhető $x(n)$ referenciajel és $e(n)$ hibajel alapján végzünk.

\feladat{5.1}

\subsubsection*{Feladat leírása}

Állítsa össze az adaptív visszhangcsökkentő rendszert, a potenciométert állítsa el a teljes kioltást megvalósító állásából! Tesztelje az algoritmus működését!

\subsubsection*{Feladat megoldása}

A DSP kártyára a \texttt{SHARC\_Lab\_LMS} projektet töltöttük le. Az $y(n)$ bemeneti jelet a kártya bemenetének bal (piros) csatornájára (a programban egy \texttt{Y\_DELAY\_BUFF\_SIZE = 80} mintás késleltető pufferen keresztül), az $x(n)$ referenciajelet a jobb (fehér) csatornájára kötöttük. A kimenetek kiosztása megegyezik a 4. feladatban leírtakkal.

\todo{A mérési elrendezés és a jelforrások leírása, a mérési eredmények.}

% \kep{../../Kepek/XXXXXXXX_XXXXXX.jpg}{A visszhangcsökkentés mérési elrendezése.}{meresi-elrendezes-lms}

\feladat{5.2}

\subsubsection*{Feladat leírása}

Próbáljon ki különféle $\mu$ értékeket, és vizsgálja a rendszer beállási és stabilitási tulajdonságait! Milyen hatással van $\mu$ változtatása a hasznos jelre?

\subsubsection*{Feladat megoldása}

\todo{Mérési eredmények. A bátorsági tényező a \texttt{processLMS.c} fájlban a \texttt{mu} változóval állítható (alapértelmezetten 0,01), a szűrő hossza $N = 100$.}

\newpage

%==========================================================================
\feladat{6}

\textit{(kiegészítő mérési feladat)} A hatodik feladatban az aktív zajcsökkentést vizsgáltuk. A DSP kártyára a \texttt{SHARC\_Lab\_ANC} projektet töltöttük le, ami identifikációs és zajcsökkentési üzemmódban működhet. Identifikációs módban (\texttt{FLAG1 / SW1} nyomva tartása) a másodlagos út identifikációja történik a kártya által generált álvéletlen zaj gerjesztéssel, zajcsökkentési módban pedig az FxLMS algoritmus fut. A \texttt{FLAG2 / SW2} gombbal a $\hat{W}(z)$ (\texttt{coef\_w}), a \texttt{DAI\_P19 / SW3} gombbal a $C(z)$ (\texttt{coef\_s}) együtthatói nullázhatók, míg a \texttt{DAI\_P20 / SW4} gombbal a kimenet tiltható (ekkor $\hat{W}(z)$ adaptálása is szünetel). A be- és kimenetek kiosztása a következő:

\begin{itemize}
    \item 1. bemenet (fehér): referenciajel ($x$ ANC módban; identifikációs módban nincs használatban),
    \item 2. bemenet (piros): hibaérzékelő jele ($e$ ANC módban; $y$ identifikációs módban),
    \item 1. és 3. kimenet: beavatkozó jel ($-\hat{y}$ ANC módban; fehér zaj identifikációs módban),
    \item 2. és 4. kimenet: hibajel ($e$).
\end{itemize}

A zavarjelet a hibrid NOISE IN bemenetére, a beavatkozó jelet az IN bemenetre adtuk, a hibajelet pedig a hibrid kiegyenlítésének megszüntetése után az OUT kimeneten érzékeltük.

\feladat{6.1}

\subsubsection*{Feladat leírása}

Identifikálja a másodlagos utat! A blokkvázlat alapján határozza meg, milyen elemek vannak a jelútban, és az identifikáció eredményét vesse össze ezzel! (Használja a memória grafikus megjelenítését!)

\subsubsection*{Feladat megoldása}

\todo{Mérési eredmények. A másodlagos út modelljének együtthatói a \texttt{coef\_s} tömbben ($I = 200$ elem) találhatók.}

\feladat{6.2}

\subsubsection*{Feladat leírása}

Periodikus zavarjelet alkalmazva tesztelje az algoritmust!

\subsubsection*{Feladat megoldása}

\todo{Mérési eredmények.}

\feladat{6.3}

\subsubsection*{Feladat leírása}

Milyen hatása van az elnyomás mértékére, illetve a beállási, stabilitási tulajdonságokra $W(z)$, illetve $C(z)$ fokszáma, valamint $\mu$ értéke változtatásának?

\subsubsection*{Feladat megoldása}

\todo{Mérési eredmények. A projektben $W(z)$ hossza $N = 100$, $C(z)$ hossza $I = 200$, a bátorsági tényezők \texttt{mu} $= 0{,}0001$ (zajcsökkentés) és \texttt{mu\_i} $= 0{,}005$ (identifikáció).}

\newpage
\section*{Függelék}
\subsection*{MATLAB kódok}

\lstinputlisting[
    caption={Az 1.1 feladat MATLAB-kódja (LMS algoritmus)},
    label={lst:lms-matlab}
]{../../MATLAB/lms.m}

\lstinputlisting[
    caption={Az 1.2 feladat MATLAB-kódja (IIR szűrő tervezése)},
    label={lst:iir-matlab}
]{../../MATLAB/tervezIIR.m}

\lstinputlisting[
    caption={Az 1.2-1.3 feladat MATLAB-kódja},
    label={lst:meres1-23-matlab}
]{../../MATLAB/Meres1_23.m}

\lstinputlisting[
    caption={Az 1.4 feladat MATLAB-kódja},
    label={lst:meres1-4-matlab}
]{../../MATLAB/Meres1_4.m}

\lstinputlisting[
    caption={A 2.1 feladat MATLAB-kódja (NLMS algoritmus)},
    label={lst:nlms-matlab}
]{../../MATLAB/nlms.m}

\lstinputlisting[
    caption={A 2.2 feladat MATLAB-kódja},
    label={lst:meres2-2-matlab}
]{../../MATLAB/Meres2_2.m}

\lstinputlisting[
    caption={A 3.1 feladat MATLAB-kódja (FxLMS algoritmus)},
    label={lst:fxlms-matlab}
]{../../MATLAB/fxlms.m}

\lstinputlisting[
    caption={A 3.2-3.4 feladat MATLAB-kódja},
    label={lst:meres3-234-matlab}
]{../../MATLAB/Meres3_234.m}

\end{document}
